% New CH-chain continuation; adopted readers and previous intakes are frozen.
\subsubsection{The finite-partial-map chain condition and the source's no-epi lemma}

The original Nice 1970 article, published in 1971, leaves the word
\emph{suitable} in its countable-chain-condition lemma on p.~333
\cite[p.~333]{Lawvere1970}. Here is an explicit sufficient internal
condition and its proof for the very finite-partial-map object used there.
These are independent editorial proofs, not unprinted formulas attributed
to the original source. The base is the declared abstract-set topos
\(\mathcal S\), with classical logic, choice, power objects and the source's
recursion object \(\omega\), not an external universe of ZFC sets.

\begin{hbdefinition}{def:p10-lawvere-ch-chain-poset}{Finite conditions and internally countable antichains}
Let \(D=\omega\times I\). Write \(\operatorname{Fin}(D)\) for its
finite-subset object, the least
subobject of \(\mathcal P D\) containing empty and singleton predicates
and closed under binary union, obtained by intersection of all such
subobjects. In the present classical recursion base these are exactly
subsets bijective with some initial segment \([n]\subseteq\omega\).
Let \(P\) be the object of functional relations \(D\rightharpoonup2\)
whose domain lies in this finite-subset object.

Write \(q\succeq p\) when \(q\) extends \(p\), the source's \(q\ge p\).
The site category has arrows \(q\to p\) in that case, so its presheaf
restriction runs toward stronger conditions. Conditions are compatible
iff they agree on their common domain; their union is then a common
extension. An antichain is a subobject of pairwise incompatible distinct
conditions. Countable means admitting a monomorphism into \(\omega\),
including finite and empty objects. The internal ccc says that every
antichain subobject is countable. This is an internal quantified statement
on the relevant power objects, valid with parameters; it is not merely an
external assertion about global conditions.

On any such poset use the double-negation coverage: a sieve at \(p\)
covers iff for every \(q\succeq p\) there is \(r\succeq q\) in the sieve.
Let \(\widetilde P\) be its internal sheaf topos and
\(u^*=a\Delta:\mathcal S\to\widetilde P\) its constant-sheaf inverse image.
The earlier internal-presheaf and reflector constructions supply this
topos and geometric adjunction. Constant sheaves are not assumed to have
only ground sections at every stage.
\end{hbdefinition}

\begin{hblemma}{thm:p10-lawvere-ch-chain-cardinal-tools}{proved=annex:p10.lawvere-ch-chain.cardinal-tools}{The bounded cardinal tools actually used}
In \(\mathcal S\), countable unions of internally countable families
are countable, images of countable objects are countable, and
\(\omega\times\omega\cong\omega\). Every infinite object \(X\) satisfies
\(X\times\omega\cong X\).
Pairwise-disjoint subfamilies and antichains have maximal subfamilies.
The annex constructs these by bounded well-ordered histories in power
objects and ordinary natural-number recursion, not external ordinal
recursion, an unproved Zorn lemma, or Replacement.
\end{hblemma}

\begin{hbtheorem}{thm:p10-lawvere-ch-chain-ccc}{proved=annex:p10.lawvere-ch-chain.ccc}{The actual finite-partial-map poset is internally ccc}
For every \(I\in\mathcal S\), the object \(P\) above satisfies the
internal ccc. The proof constructs an uncountable delta-system from an
uncountable family of distinct finite supports, then fixes the finite
root assignment. Equal root assignments make two conditions compatible.
Finite support fibers and countable unions are proved, not assumed to
preserve cardinal size.
\end{hbtheorem}

\begin{hblemma}{thm:p10-lawvere-ch-chain-constant-sections}{proved=annex:p10.lawvere-ch-chain.constant-sections}{Constant sections, deciding covers and number objects}
For a base object \(A\), the constant presheaf is separated and its
unit into \(u^*A\) is a dense mono. Thus every section of \(u^*A\)
is locally a ground section, and unequal ground sections never become
equal. These assertions hold in every generalized parameter context.
The image \(u^*\omega\), with transported zero and successor, is the
natural-number object of \(\widetilde P\).
This classical sheaf topos is a Boolean-valued model in the source's
stronger sense: its canonical free-action covers are subobjects of
base-indexed constant sheaves, and the proved matrix disjointification
therefore supplies choice and support splitting.
No enough-points assumption or logical preservation of arbitrary
function objects is made.
\end{hblemma}

\begin{hbtheorem}{thm:p10-lawvere-ch-chain-no-epi}{proved=annex:p10.lawvere-ch-chain.no-epi}{A precise sufficient no-epi preservation lemma}
Let an internal poset \(P\) satisfy the internal ccc just defined, and
let \(X,Y\in\mathcal S\), retaining the source's hypothesis
\(Y\times\omega\cong Y\). Then
\[
 \operatorname{Epi}_{\mathcal S}(X,Y)=0
 \quad\Longrightarrow\quad
 \operatorname{Epi}_{\widetilde P}(u^*X,u^*Y)=0.
\]
Here the epi objects are the full internally parametrized predicates of
surjectivity, not the absence of one chosen global map. For each ground
input, an actual maximal deciding antichain gives a countable family of
possible output values. A purported local epi yields a ground epi from
\(X\times\omega\); infinite \(X\) absorbs that factor. Finite \(X\)
is handled by deciding all its values on a common extension and choosing
a missing ground output. The proof does not make the invalid inference
that a surjection from \(X\times\omega\) always gives one from \(X\).
\end{hbtheorem}

\begin{hbexample}{def:p10-lawvere-ch-chain-naive-countercase}{Why the finite-input branch is necessary}
For \(X=1,Y=\omega\), the projection \(X\times\omega\to Y\) is
an isomorphism, while no map \(X\to Y\) is epi: a singleton image
cannot contain both zero and successor of zero. Also
\(Y\times\omega\cong Y\). Thus even the source's absorption
hypothesis does not justify the naive intermediate cardinal inference.
The proved finite-input deciding argument is what repairs that step.
\end{hbexample}

\begin{hbtheorem}{thm:p10-lawvere-ch-chain-generic}{proved=annex:p10.lawvere-ch-chain.generic}{Actual generic evaluation and local injection descent}
For the finite-partial-map \(P\), the domain where a condition assigns a
value at \((n,i)\) is a dense subpresheaf of \(\Delta(\omega\times I)\).
Its evaluation into \(\Delta2\), after sheafification, gives an actual map
\[
 u^*I\times u^*\omega\longrightarrow u^*2.
\]
Its transpose \(\alpha:u^*I\to2^\omega\), using the new sheaf topos's
power object and number object, is monic. The proof uses a fresh finite
coordinate to distinguish two unequal indices and proves the required
density and sheaf descent at arbitrary parameter bases.
\end{hbtheorem}

\begin{hbtheorem}{thm:p10-lawvere-ch-chain-comparisons}{proved=annex:p10.lawvere-ch-chain.comparisons}{The displayed strict comparisons at their internal scope}
Put \(C=2^\omega\) in \(\mathcal S\), and take \(I\) with a mono
\(C\to I\) and \(\operatorname{Epi}_{\mathcal S}(C,I)=0\), as in the
source. Such an \(I\) is available, for example \(I=2^C\), by the
singleton injection and Cantor's diagonal argument.
In the finite-partial-map sheaf topos the actual comparisons are
\[
             \omega<u^*C<2^\omega.
\]
Each strict inequality means an exhibited mono in that direction and
zero internally parametrized epi object in the same direction.
The middle-to-right mono is the transported-evaluation comparison, not
an assumed isomorphism \(u^*(2^\omega)\cong2^\omega\).
The zero epi object on the right follows from the generic injection and a complement
retraction, together with the proved no-epi lemma; no extra global sheaf
choice is assumed for that retraction.

This supplies the source's internal cardinal comparison and its omitted
ccc/no-epi/local-injection dependencies. It does not specify or prove
a translation to a separately chosen set-theoretic axiom system or model,
nor CH's positive consistency direction. Spectrum and the independently
owed general recursion/free-algebra constructions remain separate.
\end{hbtheorem}
